\begin{proof}
\begin{enumerate}[label=\arabic*.]
\item We first recover the base intrinsically.  An effective divisor on
$M_c$ cannot be horizontal: its first Chern class vanishes by
$H^2(M_c;\mathbb Z)=0$, whereas its nonzero effective restriction to a
general torus fibre would have positive pairing with the square of a
flat K\"ahler class.  The pole divisor of any meromorphic function is
therefore vertical.  On every regular fibre outside that divisor the
function is holomorphic and hence constant.  It factors through a
meromorphic function on the base, which is rational on
$\mathbb P^1$.  Thus the function field is exactly $\mathbb C(\lambda)$.
An anti-biholomorphism consequently induces an anti-M\"obius map of
the base.  The two distinct multiplicities $5,4$ and the singular cusp
support distinguish the three special values, so this base map fixes
$\infty,0,\lambda_*$.  As those points are real, the map is precisely
$\lambda\mapsto\bar\lambda$.

\item A holomorphic map between compact complex tori is affine: after
lifting to their vector spaces, its derivative is periodic and
therefore constant, since it descends to holomorphic functions on a
compact torus.  Applying this to an anti-holomorphic map after
conjugating its source gives the same assertion for its anti-linear
lift.  Its induced integral map on $H_1$ is constant along the
regular base.  Complex conjugation reverses the two finite meridians,
so this matrix must simultaneously conjugate $M_i$ to $M_i^{-1}$.
Solving these exact integral linear equations gives only $R$ and
$-R$ among unimodular matrices, where
\[
R=\begin{pmatrix}
1&0&0&0&0&0\\
2&0&0&0&-1&0\\
-1&0&-1&0&0&0\\
-1&0&0&-1&0&0\\
2&-1&0&0&0&0\\
-1&1&1&1&1&1
\end{pmatrix},\qquad R^2=I.
\]
Indeed the simultaneous intertwiner space over $\mathbb Q$ has
dimension one and this primitive matrix is unimodular; determinant
one in absolute value permits only the scalar factors $\pm1$.
The equations and this certificate are in
\texttt{centralizer\_torsion.py} and
\texttt{real\_structure\_check.py}.

\item The finite affine cap excludes $G=R$.  Normalize the pullback by
$t=u^{m_i}$; any base-preserving isomorphism lifts to this canonical
root cover and conjugates its deck generator to the inverse.
At its central torus the translation part $t_i$ would satisfy
\[
(I-M_i^{-1})t_i=-(v_i+Gv_i)/m_i\pmod\Lambda.
\]
Applying $\ell$, and using $\ell R=\ell$, gives
$0=-2\ell(v_i)/m_i$ in $\mathbb R/\mathbb Z$, impossible already
for $m_i=5$.  Hence $G=-R$ is necessary.

\item We now use the cusp marking, retaining the conjugation of based
meridians.  Conjugation sends
$g_0=(g_1g_2)^{-1}$ to $g_2g_1=g_2g_0^{-1}g_2^{-1}$.
Thus the matrix in a fixed local cusp frame is
$H=M_2^{-1}G$, and it anti-conjugates $M_0$.
Exact multiplication in the marking $S$ gives
\[
S^{-1}HS=
\begin{pmatrix}
1&0&0&-1&-1&0\\
0&1&0&-2&0&0\\
1&-1&-1&1&-1&-1\\
0&0&0&-1&0&0\\
0&0&0&0&-1&0\\
0&0&0&-1&1&1
\end{pmatrix}
=\begin{pmatrix}P&Q\\0&-P\end{pmatrix},\quad
P=\begin{pmatrix}1&0&0\\0&1&0\\1&-1&-1\end{pmatrix}.
\tag{1.17.1}
\]
Changing the based cusp path by a cusp loop multiplies by a power of
$M_0$ and does not change $P$.

\item We explain why an arbitrary translation cannot remove the obstruction
in (1.17.1).  In the additive torus coordinates whose first periods
are the $a_j$, the anti-linear map must have the form
$\zeta' =P\bar\zeta+t(q)$, because its values on these three real
periods already determine its linear part.  In multiplicative
coordinates it is
$z'=u(q)\bar z^{-P}$.
The source zero section limits to a smooth point of the central fibre.
An isomorphism takes it to another smooth point.  The central fibre
is irreducible, and its smooth locus is exactly the dense torus of
its toric normalization: every toric boundary divisor is a
normal-crossings self-identification.  In a chart about this image,
the translation therefore has the form $q^v$ times nonvanishing
holomorphic units, with $v\in\mathbb Z^3$.  Translating by the period
$v$ changes the selected component chart; the units act by torus
translation.  Neither changes the linear action $-P$ on the fan.

\item The primitive rays of the normal fan of that component are
\[
\{\pm v:0\ne v\in\{0,1\}^3\}.
\]
An extending torus monomial isomorphism must permute these rays.
This can be seen directly on a one-parameter arc $z=t^v$ meeting the
generic point of its boundary divisor: if the image vector lies in
the interior of a cone of dimension at least two, that divisor is
contracted, which is impossible for an isomorphism.  Complex
conjugation itself preserves the toric charts, so the same criterion
applies to the anti-monomial map above.
But $-Pe_2=(0,-1,1)^t$ is not one of the displayed rays.
This is a contradiction.  We have exhausted both possible integral
linear parts, so no anti-biholomorphism exists.
The assertion is explicitly about the chosen $S$ and its fan; it does
not identify different toroidal compactifications without a map
between them.

\end{enumerate}
\end{proof}
